\documentclass[11pt]{article}
\usepackage[T1]{fontenc}
\usepackage{amsmath}
\usepackage{booktabs}
\usepackage{geometry}
\usepackage{hyperref}
\geometry{margin=1in}
\hypersetup{hidelinks}

\title{Origins-DNA v0.3: Controlled Synthetic Benchmark}
\author{Adrian Lipa}
\date{15 July 2026}

\begin{document}
\maketitle

\begin{abstract}
We compare conventional DNA sequence representations with the experimental
orbital and atomic summary introduced in Origins-DNA v0.2.  The experiment uses
synthetic data, matched nucleotide composition, explicit null models, channel
ablations, repeated seeds, and bootstrap confidence intervals.  The result does
not validate a biological mechanism.
\end{abstract}

\section{Task definition}
Let $s\in\{A,C,G,T\}^{64}$.  Positive examples contain the declared motif
$m=\texttt{CGCG}$, whereas negative examples contain the cyclic decoy
$d=\texttt{GCGC}$.  Since $m$ and $d$ have the same multiset of nucleotides,
the inserted composition is matched while local order differs.  The dataset is
balanced with 400 examples for each of eight deterministic seeds.

\section{Representations and classifier}
The conventional representations are base frequencies, dinucleotide
frequencies, and positional one-hot encoding.  The experimental alternatives
are the full orbital summary, its phase-only ablation, and its class-only
ablation.  Every representation is evaluated using the same standardized
nearest-centroid classifier and stratified train/test indices.

For a test set with class recalls $r_0$ and $r_1$, balanced accuracy is
\begin{equation}
  \operatorname{BA}=\frac{r_0+r_1}{2}.
  \label{eq:balanced-accuracy}
\end{equation}
The reported intervals are deterministic percentile-bootstrap intervals for
the across-seed mean of Eq.~\eqref{eq:balanced-accuracy}.

\section{Controls}
The label-permutation null breaks the relation between training labels and
representations.  The sequence-shuffle null preserves exact nucleotide
composition within each sequence while destroying order.  These controls are
run with the same splits and metrics as the observed-label condition.

\section{Observed-label results}
\begin{table}[ht]
\centering
\caption{Mean balanced accuracy over eight seeds.}
\label{tab:results}
\begin{tabular}{lcc}
\toprule
Representation & Mean BA & 95\% bootstrap CI \\
\midrule
Dinucleotide frequency & 0.6354 & [0.6188, 0.6490] \\
Orbital classes only & 0.5375 & [0.5146, 0.5563] \\
Orbital full & 0.5365 & [0.5188, 0.5583] \\
Orbital phase only & 0.5354 & [0.5167, 0.5594] \\
Positional one-hot & 0.5312 & [0.5094, 0.5531] \\
Base frequency & 0.5250 & [0.5010, 0.5479] \\
\bottomrule
\end{tabular}
\end{table}

The dinucleotide baseline exceeds the full orbital summary by approximately
0.099 balanced accuracy units.  The orbital ablations remain close to the full
summary, so this experiment supplies no evidence that the current combination
of orbital channels adds useful order information for the declared task.

\section{Epistemic boundary}
Table~\ref{tab:results} is a synthetic software result.  It neither establishes
nor refutes a biological mechanism.  It does block the narrower claim that the
current fixed-size orbital summary already outperforms conventional
sequence-order features on this benchmark.  Negative and neutral evidence is
retained as part of the append-only project record.

\end{document}
